%----------------------------------------------------------------------------------------
% CV — Universität Greifswald / NFDI4Cat: Wiss. Mitarbeiter Biotechnologie & Enzymkatalyse
% Tailored: semantic data standards (OWL/RDF), ML-controlled biocatalysis,
% lab automation, closed-loop experimentation. Contact: Dr. Mark Dörr.
%----------------------------------------------------------------------------------------

\documentclass[10pt,a4paper,sans]{moderncv}

\usepackage[utf8]{inputenc}
\usepackage[T1]{fontenc}
\usepackage{lmodern}

\moderncvstyle{classic}
\moderncvcolor{green}

\usepackage[scale=0.78]{geometry}
\usepackage{ragged2e}

%----------------------------------------------------------------------------------------

\firstname{Kundai Farai}
\familyname{Sachikonye}

\address{Zeltnerstrasse 30}{90443, Nuernberg, Germany}
\mobile{(+49) 1626386344}
\email{kundai.sachikonye@bitspark.com}
\homepage{github.com/fullscreen-triangle}
\photo[70pt][0.4pt]{pictures/Bew2.jpg}

\quote{Bioinformatics / cheminformatics engineer: semantic data standards (OWL/RDF),
ontologies and DSLs, agentic and closed-loop research tooling --- applied to
enzyme catalysis and biocatalytic data.}

%----------------------------------------------------------------------------------------

\begin{document}

\makecvtitle

%----------------------------------------------------------------------------------------

\section{Profile}
\cvitem{}{
  Computational scientist and software engineer working at the intersection of
  \textbf{biocatalysis, semantic data, and research automation}. I design ontologies
  and domain-specific languages, build FAIR-aligned data tooling, and write the
  agentic infrastructure that drives closed-loop experimentation --- and I bring a
  working enzyme-kinetics / catalysis background to ground the data models in the
  chemistry they describe. Everything I build is deployed, documented, and open-source.
}

\section{Fit to NFDI4Cat}
\cvitem{Semantic data \& FAIR}{
  Design of ontologies, taxonomies and domain-specific languages; OWL / RDF
  knowledge representation; FAIR-data-aligned schemas for scientific exchange
}
\cvitem{(Bio)catalysis domain}{
  Enzyme kinetics and catalytic-mechanism modelling (e.g.\ carbonic anhydrase
  turnover, transition-state vs.\ geometric frameworks); mass-spectrometry data
}
\cvitem{Agentic \& closed-loop}{
  \texttt{purpose} --- an agent-queryable knowledge layer over scientific corpora;
  autonomous, iterative experiment / measurement pipelines
}
\cvitem{Engineering \& DevOps}{
  Python, JavaScript/TypeScript, Rust; Docker (Kubernetes); Git, agentic
  programming; relational / scientific databases
}

\section{Technical Skills}
\cvitem{Languages}{
  Python, JavaScript / TypeScript, Rust, Java, PHP; SQL
}
\cvitem{Semantic web}{
  OWL, RDF; ontology / taxonomy design; domain-specific languages and their
  compilers (lexer / parser / codegen)
}
\cvitem{ML / AI}{
  PyTorch; agentic / multi-agent systems; foundation-model adaptation (LoRA);
  RAG and reasoning pipelines; logic / rule-based reasoning
}
\cvitem{Infrastructure}{
  Docker, Kubernetes; Ray, Dask, SLURM; AWS, Firebase; large-scale ($>$100\,GB)
  pipelines; Git
}
\cvitem{Developer tooling}{
  5 VSCode extensions including a language server (LSP); reproducible CLI tooling
}

\section{Experience}
\cventry{2021--present}{Independent Computational Researcher \& Software Engineer}{Bitspark GmbH}{}{}{
  Built compilers / domain-specific languages, ontology-backed knowledge tooling,
  GPU pipelines and agentic AI systems --- all deployed and open-source. Authored a
  research corpus on enzyme catalysis, spectroscopy and (bio)catalytic data
  representation. Rust, Python / PyTorch, TypeScript.
}
\cventry{2019--2021}{Computational Science --- Bioinformatics, HPC \& Teaching}{Technische Universität München}{Munich}{}{
  Doctoral research in computational lipidomics; large-scale data analysis and ML
  pipelines at HPC scale (PyTorch, SLURM). Built and curated lipidomics databases.
  Taught at Master's level and supervised student work.
}
\cventry{2017--2018}{Data Pipelines}{Universitätsklinikum Hamburg-Eppendorf}{Hamburg}{}{
  Automated, reproducible high-throughput data-processing pipelines.
}

\section{Selected Projects}
\cvitem{purpose}{
  Rust CLI (Cargo workspace): an agent-queryable context layer over scientific
  project corpora --- the closed-loop / agentic-programming pattern the NFDI4Cat
  role calls for. \href{https://github.com/fullscreen-triangle/purpose}{github.com/fullscreen-triangle/purpose}
}
\cvitem{kwasa-kwasa}{
  A polyglot programming language with its own compiler and a suite of 5 VSCode
  extensions (incl.\ a language server); DSL / semantic-runtime design. Rust + TypeScript.
  \href{https://github.com/fullscreen-triangle/kwasa-kwasa}{github.com/fullscreen-triangle/kwasa-kwasa}
}
\cvitem{Enzyme-catalysis research}{
  First-principles treatment of enzyme turnover and catalytic mechanism
  (carbonic anhydrase), with platform-invariant mass-spectrometry data models
  --- domain grounding for biocatalytic data standards.
}
\cvitem{lavoisier}{
  Mass-spectrometry analysis framework: platform-invariant featurization and
  reproducible scientific pipelines (Python).
  \href{https://github.com/fullscreen-triangle/lavoisier}{github.com/fullscreen-triangle/lavoisier}
}

\section{Education}
\cventry{2018--2019}{MSc Computational Biology}{Universität Tübingen}{}{}{}
\cventry{2014--2017}{MSc Pharmaceutical Biotechnology}{HAW Hamburg}{}{}{}
\cventry{2011--2014}{BSc Biochemistry and Cell Biology}{Jacobs University Bremen}{}{}{}
\cvitem{}{\textit{Additional:} doctoral research, computational lipidomics, TUM (not completed).}

\section{Languages}
\cvitemwithcomment{English}{Native / Fluent (C2)}{required for the role}
\cvitemwithcomment{German}{Conversational}{resident in Germany since 2011}

\section{Work Authorization}
\cvitem{}{
  Permanent residence in Germany --- unrestricted work authorization, no
  sponsorship required.
}

\renewcommand{\listitemsymbol}{-~}

\end{document}
